% SOURCE_STATEMENT: "Problem 25. Give a sharp upper bound for the Morse index of a nonplanar central config-uration (as a critical point of the potential on the normalized configuration space). Give a sharp lower bound for the Morse index of a planar central configuration when it is viewed 7\n\nas part of the nonplanar configuration space. Is the Morse index related to the stability of the rigidly rotating periodic orbits? For example, does a linearly stable relative equilibrium necessarily arise from a minimum of the Newtonian potential? (Rick Moeckel)"
% Preparation and final self-audit: CONTRIBUTING.md prompts C and D; source/audit artifacts retained in the local physics_statements_2026_10_11 batch.
% !TeX program = lualatex
\documentclass[11pt]{article}
\usepackage[margin=25mm]{geometry}
\usepackage{fontspec}
\setmainfont{Latin Modern Roman}
\usepackage{amsmath,amssymb,mathtools}
\usepackage{unicode-math}
\setmathfont{Latin Modern Math}
\usepackage{xurl,hyperref}
\hypersetup{colorlinks=true,urlcolor=blue}
\setcounter{secnumdepth}{0}
\setlength{\parindent}{0pt}
\setlength{\parskip}{5pt}
\setlength{\emergencystretch}{3em}
\newcommand{\sref}[2]{\hyperlink{source:#1:#2}{[S#2]}}
\newcommand{\cataloguescope}[1]{\paragraph{Recorded target.}#1}
\newcommand{\R}{\mathbb R}
\newcommand{\C}{\mathbb C}
\newcommand{\N}{\mathbb N}
\newcommand{\Z}{\mathbb Z}
\newcommand{\Q}{\mathbb Q}
\newcommand{\F}{\mathbb F}
\newcommand{\eps}{\varepsilon}
\begin{document}
% UnsolvedMath ID 20002554; source catalogue code AIM-PHYSICS-0062
\setcounter{section}{20002553}
\section{Normal Hessian blocks, an exact four-body index jump, and stability obstructions}\label{20002554}
{\small UnsolvedMath ID 20002554 · Mathematical Physics}
\subsection{Definitions and mathematical statement}

\paragraph{Shared notation.}$\mathbb N=\{1,2,\ldots\}$; the standard number fields have their usual meanings. All additional conventions are stated locally.


\paragraph{Definitions and notation.}In the Newtonian $n$-body problem, positions $q_i\in\mathbb R^d$ and masses $m_i>0$ obey $m_i\ddot q_i=\nabla_{q_i}U$, where $U(q)=\sum_{i<j}m_im_j/|q_i-q_j|$ and the gravitational constant is normalized to one. Here $T=\tfrac12\sum_i m_i|\dot q_i|^2$, $E=T-U$, and a collision has $q_i=q_j$. The center of mass can be fixed at zero; moment of inertia is $I=\sum_i m_i|q_i|^2$. A central configuration has $\nabla_{q_i}U=-\lambda m_iq_i$ in center-of-mass coordinates for a common scalar $\lambda$. After fixing $I=1$, it is a critical point of $U$ on the collision-free inertia ellipsoid. Congruent configurations and scales are factored out for a finiteness count. A relative equilibrium is a rigidly rotating solution with time-independent shape. The Morse index of a constrained critical point is the number of negative Hessian directions, with multiplicity, after imposing the normalization and accounting for symmetry degeneracies. A planar configuration can be embedded in the spatial configuration space and acquire extra normal directions. Linear stability concerns the variational equations of the corresponding rotating motion.

\paragraph{Catalogue wording retained for traceability.}
Give a sharp upper bound for the Morse index of a nonplanar central configuration as a critical point of the positive Newtonian potential on normalized configuration space. Give a sharp lower bound for the index of a planar central configuration when viewed in the nonplanar configuration space. Determine the relation to stability of rigidly rotating periodic solutions; in particular, must a linearly stable relative equilibrium come from a minimum of that potential?
\paragraph{Scope and source qualification.}The source’s positive potential convention is made explicit. Planar versus ambient Hessians and Morse versus dynamical stability are distinct targets; the page-break artifact is removed without dropping any question. The catalogue title/status was paired with a machine-generated research attempt; that attempt is not peer-reviewed evidence of settlement. The statement covers the original source questions rather than only the title’s subcase.
\subsection{Short English statement}
Find sharp Morse-index bounds for central configurations and determine their relation to stability of rigid rotations.
\subsection{Sources}
\begin{itemize}
\item[\hypertarget{source:20002554:1}{S1}] ULAM.AI and the UnsolvedMath Contributors, UnsolvedMath, record 20002554 (AIM-PHYSICS-0062). Catalogue curation under CC BY 4.0; definitions and English exposition added locally.
\url{https://huggingface.co/datasets/ulamai/UnsolvedMath}.
\item[\hypertarget{source:20002554:2}{S2}] American Institute of Mathematics, Variational Methods in Celestial Mechanics; notes by Kuo-Chang Chen, version 22 June 2003, source item 25.
\url{https://aimath.org/WWN/varcelest/varcelest.pdf}.
\end{itemize}
\label{end:20002554}
\end{document}
